\section{Jagttermer del 3 (R-V)}

\begin{frame}
  \vspace{3em}
  \begin{center}
    {\LARGE \textbf{Jagttermer - Del 3}
    }
    \\
    \vspace{1em}
  \end{center}
\end{frame}

\begin{frame}{Jagttermer - Del 3 (R-V)}

\vspace{0.3em} % Mindre top-afstand

\begin{columns}[T]
    % --- Venstre kolonne: R, S ---
    \column{0.48\textwidth}
    \textbf{\large R, S}
    \setlength{\parskip}{0.1em} % Mindre afstand mellem punkter
    \scriptsize % Mindre skriftstørrelse
    \RaggedRight
    \begin{itemize}
        \item \textbf{Remise:} Småbeplantning i agerlandet, anlagt for vildtets skyld.
        \item \textbf{Revir:} Jagtterræn.
        \item \textbf{Rudel:} Flok af større hjortevildt (dog ikke råvildt).
        \item \textbf{Selektiv afskydning:} Målrettet afskydning af udvalgte enkeltdyr – fx dyr i dårlig fysisk stand, eller bestemt køns- og/eller aldersgruppe.
        \item \textbf{Skudtegn:} Den måde et stykke vildt reagerer på, når det træffes af skuddet.
        \item \textbf{Skydetårn:} Platform til riffeljagt eller buejagt.
        \item \textbf{Skydestige:} Stigekonstruktion med stol til riffeljagt eller buejagt.
        \item \textbf{Skørt:} Hårdusk omkring råens kønsåbning.
        \item \textbf{Smaldyr (Smalrå/smaldå/smalhind):} Hundyr hos hjortevildtarterne i tiden fra det er udvokset, til det sætter lam/kalv første gang.
        \item \textbf{Spidshjort:} Kronhjort/dåhjort/sikahjort med gevir, der kun består af to spidse stænger uden sprosser. Typisk hjortens første gevir.
    \end{itemize}

    % --- Højre kolonne: S, T, V ---
    \column{0.48\textwidth}
    \textbf{\large S, T, V}
    \setlength{\parskip}{0.1em} % Mindre afstand mellem punkter
    \scriptsize % Mindre skriftstørrelse
    \RaggedRight
    \begin{itemize}
        \item \textbf{Stand:} Særlig adfærd hos stående hunde, som på den måde angiver, at der er vildt i terrænet.
        \item \textbf{Standvildt:} Vildtarter, der forbliver i et bestemt område hele livet igennem.
        \item \textbf{Såt:} Afgrænset område, der afdrives samlet i forbindelse med en fællesjagt.
        \item \textbf{Trækvildt:} Vildt der er i en del af året trækker ud af landet, typisk fugle.
        \item \textbf{Veksel:} Hyppigt benyttet vildtsti.
        \item \textbf{Vildt:} Alle arter af fritlevende pattedyr og fugle. Definitionen fremgår i Lov om jagt og vildtforvaltning §2 (Jagtloven).
        \item \textbf{Vildtparade:} Afsluttende samling efter fællesjagt, hvor dagens nedlagte vildt lægges frem til og æres af jagtens deltagere.
    \end{itemize}
\end{columns}

\end{frame}